% ------------------------------------------------------------------------------
% Chapter 3
% ------------------------------------------------------------------------------
\chapter{Diagram} 
\label{cha:diagram}

For simplicity, the main diagram (figure \ref{fig:diagram}) has the main structural pieces separated, while they follow a hierarchical structure.
That structure is shown in figure \ref{fig:hierarchy}.

\begin{figure}[htb]
  \centering
  \includesvg[width=\textwidth, inkscapelatex=false]{../diagram.drawio.svg}
  \caption{Diagram of the structure of the testbench}
  \label{fig:diagram}
\end{figure}

There are multiple sequences to improve the modularity of the system: random sequence and deterministic write + read.

The green arrows represent the communication that is handled by the UVM ports.
The subscriber uses the green rectangle to make the "man in the middle" attack to that communication.

Each interface is connected to its own driver and the pins to a single master in the \texttt{DUT}.
This improves modularity, but creates a second problem: the concurrence information is lost in this communication.
The scoreboard can not see if 2 consecutive transactions have been executed together.

Both the \texttt{checker} and the \texttt{coverage} modules are connected to the \texttt{DUT} through a binding.
This connection means that they are not actually connected to the interface, but incorporated to the module itself, but for simplicity they were connected to the interface here.
These two modules are in charge of the concurrent behavior of the two masters connected to the \texttt{DUT}.


\begin{enumerate}
    \item \texttt{Testbench}: This component simply wraps around every component and build the final machine.
    \item \texttt{Top}: This component instantiates the \texttt{DUT}, its interfaces the test to be executed and the non UVM modules such as checker and cover.
    \item \texttt{Test}: This component is responsible for the creation of sequences and the very environment that will perform all UVM transactions.
    \item \texttt{Environment}: This component is the core of the UVM structure. Here there are the agents (sequencer, driver, and monitor), and the scoreboard. 
\end{enumerate}

\begin{figure}[h!]
  \centering
  \includesvg{../hierarchy.svg}
  \caption{Hierarchy of the testbench components}
  \label{fig:hierarchy}
\end{figure}